\section{The Silicon Labs Toolchain, From the Outside}
\label{sec:vendor-toolchain}

Everything in this book is built with open tools: \codeinline{cargo}, \codeinline{rustc},
\codeinline{probe-rs}, \codeinline{openocd}, \codeinline{llvm-objdump}. And yet every binary in it
would be impossible without a vendor: the register definitions come from Silicon Labs' SVD, and the
radio stacks are their binary blobs. The vendor's tools, therefore, are not optional — but most of
them are not needed either, and knowing which is which saves weeks.

This section is about that boundary: what Silicon Labs ships, what each piece is actually for, and
what this project measured about the three that touch a Rust workflow.

\subsection{The Pieces, and What They Are For}

\begin{table}[htbp]
    \centering
    \footnotesize
    \begin{tabular}{@{}p{0.24\textwidth}p{0.36\textwidth}p{0.32\textwidth}@{}}
        \toprule
        \textbf{Piece} & \textbf{What it is} & \textbf{Does a Rust path need it?} \\
        \midrule
        Simplicity Studio 5 & IDE, part database, board files, configurators
            & No for building; \textbf{yes} for the Radio Configurator \\
        Gecko SDK (now Simplicity SDK) & headers, emlib C sources, prebuilt stack libraries
            & \textbf{Yes}: SVD, and the blobs \\
        Simplicity Commander & vendor CLI for flashing, locking, reading flash
            & No, and it is actively dangerous here (see below) \\
        SEGGER J-Link & external debug probe, licensed
            & No: the board has its own debugger \\
        CMSIS-DAP on the board & the SAMD11 chip that \emph{is} the debugger
            & \textbf{Yes}, via \codeinline{probe-rs} or \codeinline{openocd} \\
        Gecko Bootloader & the factory image at the base of flash
            & \textbf{Yes}, and it constrains every linker script \\
        AppLoader UART API & OTA over a serial line
            & Only for OTA; it is a protocol, not a tool \\
        Radio stack blobs & \codeinline{libbluetooth_*}, \codeinline{librail_*}
            & \textbf{Yes}, under MSLA licensing \\
        \bottomrule
    \end{tabular}
    \caption{The vendor surface. Two things are unavoidable — the SVD and the radio blobs — and
    everything else is a convenience or a hazard.}
    \label{tab:vendor-pieces}
\end{table}

\subsection{JTAG, SWD, and the Debugger Already on the Board}

There are three ways to reach a Cortex-M core, and the naming confuses two of them:

\begin{itemize}
    \item \textbf{JTAG} — the five-wire standard (TCK, TMS, TDI, TDO, nTRST). Everything supports it,
          nothing needs it on a modern single core.
    \item \textbf{SWD} — two wires (SWCLK, SWDIO). It is what the EFR32MG24 uses, and what every
          board-level debug connector exposes.
    \item \textbf{cJTAG} — two-wire \emph{JTAG}, a different protocol that happens to use two pins.
          Some Silabs documentation mentions it; it is not what \codeinline{probe-rs} speaks.
\end{itemize}

The XIAO MG24 Sense has no debug header. It has a second microcontroller — a SAMD11 — wired between
the USB connector and the EFR32's SWD pins, and that chip implements the CMSIS-DAP protocol
(USB \codeinline{2886:0062}). Two consequences, both measured:

\begin{itemize}
    \item \codeinline{probe-rs} sees the EFR32 through that bridge and can program it. So can
          \codeinline{openocd} with \codeinline{interface/cmsis-dap.cfg}. Neither needs a J-Link.
    \item The bridge is \emph{flaky} in a specific way: the first attach works, and a later one dies
          with \codeinline{Command ID in response (0x14) does not match sent command ID}. A USB reset
          does not clear it; unplugging the cable does. That is a fact about the SAMD11, not about
          your program, and it will waste an afternoon if you assume otherwise.
\end{itemize}

\subsection{The Bootloader Is Not Optional, and It Is Why Linker Scripts Matter}

Every XIAO MG24 Sense leaves the factory with a Gecko Bootloader at the base of flash. It occupies
the first 24 KiB, and on this part the application therefore starts at \textbf{0x08006000}, not at
\codeinline{0x08000000}.

That single fact decides the linker script, and getting it wrong does not produce an error message —
it produces a board that no longer boots, because the image was written over the bootloader. The
symptom is subtle enough to be worth stating in full: the write succeeds, the read-back matches, and
the board is silent.

Hence the two requirements this project imposes on itself, both mechanical:

\begin{minted}{text}
/* 1. never let the linker put anything below the bootloader (memory.x) */
FLASH (rx) : ORIGIN = 0x08006000, LENGTH = 1536K - 24K
\end{minted}

\begin{minted}{bash}
# 2. refuse to flash an image that would
python3 scripts/elf-load-range.py "$ELF"     # load_min must be >= 0x08006000
\end{minted}

\subsection{Simplicity Commander: Convenient, and the Wrong Tool Here}

\codeinline{commander} is the vendor's CLI: it flashes, reads flash, locks and unlocks parts, and
writes the bootloader. It is well documented and it does not check where your image says it goes. A
\codeinline{commander flash image.hex} writes whatever the hex file describes — including, if the
linker script is wrong, 0x08000000.

For a workflow that must not brick a board, that is the wrong property. The alternative this project
uses is a script that (1) inspects the ELF's load range, (2) refuses anything starting below
0x08006000, (3) writes with \codeinline{probe-rs download}, and (4) verifies with
\codeinline{probe-rs verify}. The guard is the point; the tool is replaceable.

\subsection{The Radio Stack: Two Licences in One Directory}

This is the part of the vendor surface that a Rust project cannot avoid and cannot reimplement.

\begin{itemize}
    \item \textbf{Headers and emlib sources} are Zlib licensed: usable, redistributable, and
          readable. \codeinline{sl_bt_api.h} is where the BGAPI command set is defined, and
          \codeinline{rail.h} is where RAIL is.
    \item \textbf{The stacks themselves} are binary: \codeinline{libbluetooth_host},
          \codeinline{libbluetooth_controller}, \codeinline{libbgcommon} for Bluetooth, and
          \codeinline{librail} for the proprietary/802.15.4 side. They are MSLA licensed, which is
          why they are not in this book's repository and why a clean checkout does not link.
\end{itemize}

Two traps follow from that split, and both cost real time in this project:

\textbf{The header and the blob must be the same SDK version.} When they are not, the blob exports
symbols the bindings do not declare — \codeinline{sl_bt_advertiser_create_set} was in the blob and
absent from the generated bindings — and the failure looks like a typo in your code. It is not.

\textbf{A blob references code that is not in the blob.} The Bluetooth stack calls the CMSIS
\codeinline{CORE_*} helpers; RAIL calls four emlib functions
(\codeinline{CMU_ClockEnable}, \codeinline{EMU_DCDCSetPFMXTimeoutMaxCtrl},
\codeinline{SYSTEM_ChipRevisionGet}, \codeinline{SystemSYSCLKGet}). The application is expected to
provide them, because in a vendor project the linker pulls them from emlib. In a Rust project they
have to be provided deliberately — and the CMSIS ones must be provided \emph{once}, for both stacks,
because two definitions of \codeinline{CORE_EnterCritical} either collide or, worse, disagree about
which interrupt mask they set.

\subsection{The Radio Configurator, and Why RAIL Needs a Generated File}

RAIL is a library, not a radio. Before it can transmit it needs a \emph{radio configuration}:
channel numbers and their modem settings, the PA and front-end tables for this board's matching
network, and the calibration data for the crystals actually fitted. Nothing in the header contains
those numbers, and no amount of reading the reference manual produces them.

They are generated. The Radio Configurator inside Simplicity Studio emits \codeinline{rail_config.c}
from a description of the board and the desired protocol, and that file is the difference between
"the library initialised" and "a packet left the antenna". This is the one place where the vendor
IDE is genuinely required, and it is worth saying plainly: a Rust project can avoid Studio for
editing, building, flashing and debugging, and cannot avoid it for radio.

\subsection{What To Take Away}

\begin{itemize}
    \item The vendor tools split into: \emph{data you need} (SVD, headers, emlib, blobs), \emph{a
          file you cannot write yourself} (the radio configuration), and \emph{conveniences that
          fight your safety rules} (Commander).
    \item The board's own debugger is a CMSIS-DAP bridge; SWD is the two-wire protocol; J-Link is
          unnecessary, and the bridge needs a physical replug when it desynchronises.
    \item The bootloader owns the bottom of flash. A linker script that ignores it bricks the board
          silently, so the guard belongs in a script that runs before the write.
    \item Zlib headers plus MSLA blobs is a normal arrangement, and it has two failure modes worth
          remembering: version mismatch, and symbols the blob expects the application to supply.
\end{itemize}
